\section{Coding Agent Index v2.1、Claude系が首位}
Artificial AnalysisがCoding Agent Index v2.1を公開した。FrontierCodeを2.0、Terminal-Benchを4.1に更新し、エージェント実行環境でのコーディング能力を測る。Grok原文によればClaude Fable 5.2とClaude Opus 5.5が58点で同率首位となった。評価更新で平均スコアが下がる一方、実行コストには約30倍の開きがあると報告された。性能順位だけでなく、harnessと費用を含めた比較が重要になっている。
\dailyxsource{Artificial Analysis公式投稿}{https://x.com/ArtificialAnlys/status/2121393821955334328}

\section{DeepSeek、検索・OCR・音声をAPI提供との報告}
DeepSeekがV4.1 Flash向けに、引用付きWeb検索、ファイル・URL処理、OCR、音声認識と音声合成などの機能をサーバー側APIとして提供したとの投稿が広がった。アプリに組み込まれていた機能を開発者が直接呼び出せるという説明で、スタンドアロンAPIにも言及がある。Grokの判定はLikely。公式発表は11月13日とされ観測窓外であり、今回の対象は窓内の利用者による紹介・反応である。
\dailyxsource{窓内の紹介投稿}{https://x.com/corbintips/status/2121722051471745140}

\section{Qwen GGUF向け量子化改善の主張}
NaceAIがAIエージェントを使ってQwenモデルのテンソル精度配分を調整し、GGUF量子化の品質を改善したとの投稿があった。一部のモデルではnext-token agreementが90％を超えるという主張もある。容量と精度の両立はローカル実行の実用性に直結するが、Grokは公式一次資料を確認できておらずUnverifiedと判定した。数値を確定的な評価結果として扱わず、公開物と再現可能な比較の確認が必要だ。
\dailyxsource{関連投稿}{https://x.com/TeksEdge/status/2121724857343705176}

\section{Antigravity /teamworkの研究利用に関心}
Google ResearchのAntigravity /teamworkについて、研究者と複数のAIエージェントが協調し、数学や化学の未解決課題に成果を挙げたとの紹介が続いた。情報理論の予想の証明などが例として挙げられ、外部利用が可能になったとの主張も共有された。Grokの判定はLikelyで、公式発表は観測窓直前の11月14日05時54分頃とされる。本号は窓内の反応を記録し、成果の範囲や利用条件の独立確認は行っていない。
\dailyxsource{窓内の関連投稿}{https://x.com/ZST_WIP/status/2121724878021922882}

\section{OpenAgentAuth、エージェント認証標準の提案}
AIエージェント向けのオープン標準「OpenAgentAuth」を公開したとする投稿があった。Ed25519を用いる識別、権限委任、サーキットブレーカーなどを備えると説明され、npxで利用するコマンドも示された。自律的に外部操作を行うエージェントの権限をどう制御するかは重要な課題だが、Grokは公式リポジトリを確認できずUnverifiedとした。標準としての採用状況や実装内容は未確定である。
\dailyxsource{開発者の投稿}{https://x.com/AmanNanda/status/2121421804762010043}

\section{AgentOne、ローカル対応デスクトップエージェント}
オープンソースのデスクトップAIエージェント「AgentOne」が紹介された。投稿によればMCPに対応し、ローカルモデルも利用でき、ライセンスはAGPL-3.0という。外部サービスに依存しないエージェント実行環境を求める利用者には関心のある方向性だ。GrokはLikelyと評価したが、投稿で案内された返信先のリポジトリは確認できていない。具体的な機能や対応OS、保守状況については追加情報を待つ必要がある。
\dailyxsource{開発者の投稿}{https://x.com/DanKornas/status/2121642335955992627}

\section{Gemini 4の利用枠とアクセス制限を巡る議論}
Gemini 4 Hydrogenの無料枠やレート制限、Carbonの利用可能性を巡ってユーザー投稿が続いた。高速モデルを実際の開発環境でどの程度使えるかが関心の中心で、利用報告の間には差も見られる。公式発表は観測窓外であり、Grokは今回の具体的な制限情報をUnverifiedと判定している。各利用者の経験をサービス全体の正式な仕様と混同せず、提供プランや地域による違いに注意したい。
\dailyxsource{利用者投稿}{https://x.com/buildwithhassan/status/2121724502220456054}

\section{Grok Botのharness最適化に関する数値}
Grok Botのharnessを最適化した結果、同一プランで利用量が約7％、ルーティン処理では最大15％増えたとの投稿が共有された。モデル自体を変えずに実行制御やツール呼び出しを改善する考え方は、エージェント運用の費用対効果を左右する。ただしGrokの判定はUnverifiedであり、xAIまたはGrokの公式発表は確認されていない。数値は利用者が共有した主張として扱う。
\dailyxsource{関連投稿}{https://x.com/MichaelGannotti/status/2121712617689383184}

\section{Runway Praxis-1、低価格ロボットアーム対応の噂}
RunwayのロボットモデルPraxis-1が低コストのSO-101アームに対応し、RTX 5090単体で実行でき、チェックポイントも公開されたとの見出しが共有された。ロボット学習モデルの実行環境を小型化できれば実験の敷居は下がるが、今回の情報はまとめ投稿に基づく。Grokは公式一次資料を確認できずUnverifiedとした。対応ハードウェア、必要な計算資源、公開条件はいずれも未確認である。
\dailyxsource{関連投稿}{https://x.com/DadWhoBuilds/status/2121723903835873589}

\section{コーディングエージェントのharnessと費用差}
「エージェントはモデルとharnessの組み合わせ」という観点から、コーディングエージェントの実行環境設計とコストを論じる投稿があった。モデルの性能が同等でも、タスク分割、ツールの使い方、再試行の仕組みによって実際の結果が変わるという問題意識だ。Artificial Analysisの評価で示された大きなコスト差とも関連づけられた。投稿ではWavestoneの論文にも触れるが、詳細はGrok収集段階では未確認である。
\dailyxsource{分析投稿}{https://x.com/buildAgoat/status/2121712174758543389}
